\section{Commercial missing-data workflows and transfer between bank populations}
\label{exec21:transport-scope}

A bank may obtain a credible intervention estimate in a pilot and still make a
poor decision when it extends the intervention to another branch. Customers may
differ, follow-up outcomes may be missing, and the software used to complete
the records may report uncertainty for a different quantity from the proposed
decision. A useful commercial comparison therefore follows the complete
question: which population receives which action, when is the outcome measured,
why is it unobserved, and which assumptions connect the observed records to the
decision? Imputation, uncertainty combination and population transfer are
separate steps. A feature called ``missing values'' establishes none of those
connections by its name alone.

The selected primary documentation provides concrete comparators. Stata's
\texttt{mi} family manages multiple completed datasets and their analysis;
SAS/STAT's MI and MIANALYZE procedures separate generation and combination;
IBM SPSS Missing Values exposes imputation and output-specific pooling.
The inspected IBM edition is 31.0.2. The SAS HTML identifies version 14.1,
whereas the Stata manuals contain release-19 example-data paths and July 2026
PDF metadata. These identities delimit the comparison; later retrieval does
not authenticate every byte as the edition available at the literature cutoff.
Exact licensed builds, current entitlements and broader commercial coverage
remain unresolved. EViews, Mplus and commercial graphical tools remain leads
requiring matched technical originals for this particular target.
\cite{exec21_transport_stata_mi,exec21_transport_sas_mi_intro,exec21_transport_spss_missing}

\subsection{Missingness assumptions precede record completion}
\label{exec21:transport-missingness}

Let $A$ be a binary action, $X$ covariates recorded before that action, $Y(a)$
the monetary outcome under action $a$, and $R$ an indicator that the outcome is
observed. The record contains $(X,A,R,RY)$, with a separate missingness mask so
that an absent outcome is never mistaken for a zero. Consistency connects an
observed outcome to $Y(A)$; treatment exchangeability and positive treatment
probability are additional causal assumptions. For the particular problem in
which only $Y$ is missing, a conditional mean restriction sufficient for response
adjustment is
\begin{equation}
 E[Y\mid X,A,R=1]=E[Y\mid X,A],\qquad
 P(R=1\mid X,A)>0.
 \label{exec21:transport-mar}
\end{equation}
The stronger familiar MAR restriction makes response independent of the missing
outcome conditional on the observed variables. When covariates are also missing,
the pattern-specific observed information and imputation model must be specified
again; equation~\eqref{exec21:transport-mar} is not a general multivariate recipe.
Administrative censoring, a variable undefined for a customer, and an outcome
lost through nonresponse need distinct observation models.

Stata's substantive manual and the SAS assumptions chapter distinguish MAR
from MCAR and describe ignorability using MAR together with distinct parameters
for the data and missingness models. They also explain why an assumption about
unobserved values cannot be verified by inspecting only the observed values.
The SAS profile adds distributional restrictions for particular regression and
MCMC procedures. These restrictions must accompany a procedure name. A good
fit to observed records, a plausible-looking imputed distribution or rejection
of an MCAR diagnostic does not establish MAR.
\cite{exec21_transport_stata_substantive,exec21_transport_sas_assumptions}

An exact example makes the investment consequence visible. Let a latent response
type $H$ be equally likely to be zero or one, randomize $A$ independently of
$H$, and set $R=H$. All customers have $Y(0)=0$. Responders have $Y(1)=10$.
In one complete-data law, nonresponders have $Y(1)=-10$; in a second they have
$Y(1)=10$. The observable joint distribution of $(A,R,RY)$ is identical in
both laws: each action-response cell has probability $1/4$, and observed treated
outcomes equal 10. Yet the complete-population effects are respectively 0 and
10. At an assumed action cost of USD 4, the margins are USD $-4$ and USD 6.
More imputation draws from one selected model cannot distinguish these worlds.
The project reference enumerates the complete laws and verifies equality of
their observed distributions. It does not assert that either law describes
actual bank records.\cite{exec21_transport_reference}

A transparent sensitivity parameter writes the unobserved treated mean as
$10+\delta$, keeping the control mean zero. The resulting target is
\begin{equation}
 \tau(\delta)=10+\tfrac12\delta,\qquad
 V(\delta)=6+\tfrac12\delta.
 \label{exec21:transport-sensitivity}
\end{equation}
For the illustrative restriction $\delta\in[-20,0]$, value ranges from $-4$
to 6 and changes sign at $\delta=-12$. The interval is supplied as a sensitivity
restriction, not learned from the missing outcomes. The SAS sensitivity chapter
documents a pattern-mixture approach in MI and explicitly says that the
selection-model approach is not supplied by that procedure profile. A result
that changes under a plausible MNAR scenario establishes dependence of the
decision on assumptions; it does not establish which scenario is true. Exact
syntax and later-release capabilities need their own evidence.
\cite{exec21_transport_sas_sensitivity,exec21_transport_reference}

\subsection{What the three documented workflows preserve}
\label{exec21:transport-workflows}

Stata offers four MI storage styles while retaining the original data as
$m=0$ and the completed datasets as $m=1,\ldots,M$. Variables can be regular,
imputed or passive, with passive variables representing algebraic functions
of other variables. The introduction also states that extended missing values
\texttt{.a} through \texttt{.z} persist as hard missing in the documented
profile. An interchange layer that maps every absent value to an ordinary
numeric null could therefore change what the engine is asked to impute.
Maintaining the missingness reason and role is part of the statistical contract.
\cite{exec21_transport_stata_mi}

For chained equations, Stata documents a default burn-in of ten iterations,
ordering from the most to the least observed variable, configurable conditional
models and saved trace summaries. Perfect prediction produces an error unless
an explicit augmentation option is requested. The manual discusses
incompatible conditionals and dependence of results on specification, order
and burn-in. These are useful controls, but a chosen default is not a proof of
convergence or of adequacy for a causal estimand. For example, an imputer that
omits a treatment interaction needed by the analyst can erase the heterogeneity
needed for the branch-transfer decision even if marginal summaries look
reasonable.\cite{exec21_transport_stata_chained,exec21_transport_stata_substantive}

There is a stronger reason to retain the conditional specifications themselves.
Consider binary $U,V$ with proposed conditionals
\begin{equation}
 \begin{aligned}
 P(U=1\mid V=0)&=P(U=1\mid V=1)=\tfrac12,\\
 P(V=1\mid U=0)&=\tfrac14,\qquad P(V=1\mid U=1)=\tfrac34.
 \end{aligned}
 \label{exec21:transport-conditionals}
\end{equation}
For any positive binary joint distribution, the odds ratio calculated from
either pair of conditionals must agree. Here the first specification requires
odds ratio 1 and the second requires 9, so no such joint distribution exists.
Nevertheless, alternating draws can have stable distributions. In the cell
order $(0,0),(0,1),(1,0),(1,1)$, a scan updating $U$ then $V$ has stationary
law $(3,1,1,3)/8$; reversing the scan has stationary law $(1,1,1,1)/4$.
Direct multiplication verifies each invariant distribution. Thus an apparently
stable trace does not demonstrate compatibility of the supplied conditionals.
This exact constructed example is a diagnostic design, not an observed failure
of a particular vendor's implementation.\cite{exec21_transport_reference}

SAS supplies a related but separately versioned workflow. Its MI overview
distinguishes sequential methods for monotone patterns from multivariate-normal
MCMC and fully conditional specification for arbitrary patterns. FCS uses
separate conditional models rather than beginning with an explicit common
multivariate model. The completed datasets are analyzed using the relevant
complete-data procedures, and MIANALYZE reads parameter estimates together with
standard errors or covariance matrices. It notes an additional assumption for
multivariate Wald inference and special treatment for quantities such as
correlations and ratios. The proposal therefore needs a typed collection of
estimates and covariance matrices, not just a filled table. The recovered HTML
contains equations as images that are outside this text inspection; the numerical
pooling formulas below are grounded in the readable Stata manual.
\cite{exec21_transport_sas_mi_intro,exec21_transport_sas_mianalyze_intro,exec21_transport_sas_combining}

IBM's manual makes another important distinction. Multiple Imputation accepts
analysis weights in summaries and imputation models, but it does not explicitly
handle strata, clusters or other complex sampling structures. The manual also
says that Complex Samples procedures do not automatically analyze MI datasets
in this profile. Frequency weights and analysis weights have different
semantics; noninteger frequency weights are rounded. Passing a final sampling
weight therefore does not establish a cluster-aware covariance calculation.
Stata documents a survey-prefix analysis route, but that alone also leaves
the imputation model's treatment of the sampling design to be justified.
\cite{exec21_transport_spss_missing,exec21_transport_stata_estimate}

SPSS chooses a monotone method when the pattern permits it and otherwise uses
FCS. Automatic FCS uses ten iterations; the manual provides history output
for convergence assessment. Main effects are the automatic specification,
with optional categorical two-way interactions under selected settings.
Scale variables can use linear regression or predictive mean matching, whose
documented default donor count is five; categorical models use logistic
regression. Constraints can cause redraws, and excluding highly incomplete
variables removes them from both imputation and prediction. Reproducing a run
requires the same seed, case order, variable order and settings. These details
affect the implied procedure and belong in an adapter request.
\cite{exec21_transport_spss_missing}

The output variable \texttt{Imputation\_} distinguishes original data from
the $M$ completed versions, and split-file settings control how those versions
are analyzed. IBM distinguishes naive combination, which supplies a pooled
parameter, from univariate combination, which can also supply standard errors,
effective degrees of freedom, intervals and missing-information diagnostics.
Support is specific to each output; an unsupported procedure returns no pooled
output. An integration must preserve that unavailability rather than attach a
generic confidence interval to every pooled number.
\cite{exec21_transport_spss_missing}

\subsection{Combine uncertainty on the declared estimand scale}
\label{exec21:transport-pooling}

For a fixed scalar estimand $Q$, let $\widehat Q_m$ and $U_m$ be the estimate
and its completed-data variance from imputation $m$. The documented combination
rule is
\begin{align}
 \overline Q&=M^{-1}\sum_{m=1}^M\widehat Q_m,&
 \overline U&=M^{-1}\sum_{m=1}^M U_m,\nonumber\\
 B&=(M-1)^{-1}\sum_{m=1}^M(\widehat Q_m-\overline Q)^2,&
 T&=\overline U+(1+M^{-1})B.
 \label{exec21:transport-poolformula}
\end{align}
The between-imputation term reflects missing-data uncertainty under the
specified imputation model, and $B/M$ reflects using finitely many imputations.
Neither accounts automatically for uncertainty about an unverified MAR or
population-invariance assumption. Valid completed-data inference and suitable
imputation remain premises. For vector parameters the covariance matrices,
including their off-diagonal terms and common coefficient identities, must be
preserved.\cite{exec21_transport_stata_estimate,exec21_transport_stata_substantive}

For the authored values $(\widehat Q_1,\widehat Q_2,\widehat Q_3)=(0,2,4)$
and $U_m=1$, equation~\eqref{exec21:transport-poolformula} gives
$\overline Q=2$, $B=4$ and $T=19/3$. Retaining only the within variance gives
1. Incorrectly treating the three completed-data estimates as independent
replicate measurements would give $\overline U/M=1/3$. Neither is the MI
variance. The documented large-complete-sample degrees-of-freedom approximation
uses $r=(1+M^{-1})B/\overline U$ and
$\nu=(M-1)(1+1/r)^2$, giving $361/128$ here. A small complete-data sample
requires the further adjustment documented by the manual; the arithmetic is
not a universal coverage result.\cite{exec21_transport_stata_estimate,exec21_transport_reference}

The estimand scale also matters. If two log-odds estimates correspond to odds
ratios 1 and 9, exponentiating their average coefficient gives 3, whereas
averaging the odds ratios gives 5. Stata documents coefficient-scale pooling
followed by an optional exponentiated display; it also has explicit requests
for transformed coefficients. A declared transformed estimand needs its own
completed-data uncertainty and suitable pooling justification. Neither 3 nor
5 is a bank's average profit or a marginal treatment effect merely because it
was returned after imputation.\cite{exec21_transport_stata_estimate,exec21_transport_reference}

Failure handling is part of the same argument. Stata normally stops when an
imputation's estimation fails or the estimation sample changes. Its
\texttt{errorok} and \texttt{esampvaryok} options permit different behavior.
An automatic wrapper must expose such choices and retain the failing
imputation and sample identities. Otherwise it can silently pool different
populations or select only successful fits. The appropriate number of
imputations should be chosen using stability and Monte Carlo error diagnostics
for the requested inference; high relative efficiency of a point estimator
does not guarantee stable intervals or decision thresholds.
\cite{exec21_transport_stata_estimate,exec21_transport_sas_mi_intro}

\subsection{Transport needs an explicit population bridge}
\label{exec21:transport-target}

Let $S=1$ denote the source pilot and $S=0$ the target branch, and define
$p_s(x)=P(X=x\mid S=s)$ and
$\tau(x)=E[Y(1)-Y(0)\mid X=x,S=1]$. One sufficient route assumes consistency,
source treatment exchangeability, positive source treatment and response
probabilities, valid response adjustment within source covariates, and equality
of conditional potential-outcome means across $S$. It also requires source
support wherever $p_0(x)>0$. Under these restrictions, iterated expectation gives
\begin{equation}
 \tau_0=\sum_x p_0(x)\tau(x)
       =E\left[\left.\frac{p_0(X)}{p_1(X)}\tau(X)\right|S=1\right].
 \label{exec21:transport-standardize}
\end{equation}
This is a project derivation for a declared finite-covariate regime. Equality of
conditional means is extra scientific information; it is not established by
the existence of a weight column or by the imputation procedure.
\cite{exec21_transport_reference}

For clarity, write $e_a(x)=P(A=a\mid X=x,S=1)$ and
$r_a(x)=P(R=1\mid A=a,X=x,S=1)$. Under the same response restriction and
positivity, each source conditional mean can be written
\begin{equation}
 \mu_a(x)=E\left[\left.
       \frac{\mathbb{1}\{A=a\}RY}{e_a(x)r_a(x)}
       \right|X=x,S=1\right].
 \label{exec21:transport-response}
\end{equation}
Condition first on $A=a$ and $R=1$, cancel the treatment and response
probabilities, then apply the mean restriction and source causal assumptions.
Substitution into equation~\eqref{exec21:transport-standardize} gives a
combined population, treatment and response-weight expression. It is not a
variance estimator: estimated nuisance models, sampling the target covariates,
clustering and any use of MI require a jointly justified inference procedure.
Weights for these different purposes must remain separately named.

Suppose the source proportions of two customer strata are $(4/5,1/5)$ and the
target proportions are $(1/5,4/5)$. Conditional intervention effects are USD
10 and USD $-10$. The pilot's average effect is 6, the target's is $-6$, and
the population weights are $1/4$ and 4. At an assumed cost of USD 2, the pilot
margin is USD 4 and the target margin is USD $-8$. The exact reference also
enumerates randomized treatment and four positive response rates, verifying
that the combined weighting expression recovers the stipulated target.
No software error is needed for the naive deployment decision to reverse:
the wrong population suffices.\cite{exec21_transport_reference}

If the source instead contains only the first stratum, the second effect is
unobserved and no finite reweighting creates its information. With the explicit
restriction that the second effect lies in $[-10,10]$, the target effect is
$2+(4/5)\tau(1)\in[-6,10]$, giving margin $[-8,8]$ at cost 2. Without such
a justified bound, even that interval cannot be asserted. An adapter should
report the support defect and any defensible identified set. Silently trimming
the second stratum changes the target population; silently predicting into it
adds extrapolation assumptions.\cite{exec21_transport_reference}

\subsection{Competitive position and a public qualification route}
\label{exec21:transport-integration}

The documented commercial systems already provide substantive imputation,
analysis and diagnostic workflows. The proposal has not demonstrated a native
replacement or superior estimation performance. Its proposed contribution is
to preserve the causal question and its restrictions across those workflows:
source and target populations, observation clocks, missingness reasons,
imputation identity, covariance semantics and explicit refusal. Licensed
isolated execution is a candidate integration route; public manuals do not
grant redistribution rights for the engines. EViews, Mplus and graphical-tool
profiles cannot be ranked without the missing originals, and prices remain
unknown.\cite{exec21_transport_stata_estimate,exec21_transport_sas_mianalyze_intro,exec21_transport_spss_missing}

A useful first prototype would compare a selected licensed procedure and an
already reviewed open implementation on the same generated population,
missingness law and estimand. Complete-record analysis, complete-case analysis,
single imputation, proper multiple imputation and an explicitly specified
response-weight procedure supply substantive alternatives. Fix the causal
target, sample law, interaction structure, number of imputations, diagnostic
and refusal rules before evaluation. Include MAR and MNAR cases, complete and
missing source support, compatible and incompatible imputation models, cluster
sampling, changing coefficient sets and unsuccessful fits. Report bias,
coverage under the claimed regime, interval width, unavailable requests,
simulation error and full workflow cost. A sensitivity grid evaluates
robustness; it is not a device for choosing whichever assumption makes a
preferred intervention profitable.

The library, CLI and MCP interfaces should carry the same request and result
meaning. Original and imputed data need stable unit and imputation identifiers,
measurement types, category references, passive transformations and sampling
design. Results need per-imputation estimates and covariance, pooled scale and
degrees of freedom, convergence diagnostics, population support and retained
failures. An assistant can help locate technical passages, propose models and
explain sensitivity; it cannot infer the unobserved outcomes or authorize a
scientific assumption from fluent text. These public comparisons can establish
whether the proposed integration merits a more expensive bank experiment.
